% Overleaf: Menu -> Compiler: pdfLaTeX (mặc định)
\documentclass[aspectratio=169,10pt]{beamer}
\usetheme{Madrid}
\usecolortheme{default}
\setbeamertemplate{navigation symbols}{}

\usepackage[utf8]{inputenc}
\usepackage[T5]{fontenc}
\usepackage[vietnamese]{babel}
\usepackage{listings}
\usepackage{xcolor}
\usepackage{booktabs}

% ---------- Dart cho listings ----------
\lstdefinelanguage{Dart}{
  keywords={abstract,as,async,await,bool,class,const,double,else,extends,
    false,final,for,if,import,int,late,new,null,return,super,this,true,
    var,void,with},
  sensitive=true,
  comment=[l]{//},
  morecomment=[s]{/*}{*/},
  morestring=[b]',
  morestring=[b]"
}
\lstset{
  language=Dart,
  basicstyle=\ttfamily\scriptsize,
  keywordstyle=\color{blue!70!black}\bfseries,
  commentstyle=\color{green!45!black},
  stringstyle=\color{red!60!black},
  backgroundcolor=\color{gray!8},
  frame=single, rulecolor=\color{gray!40},
  columns=fullflexible, keepspaces=true,
  showstringspaces=false, breaklines=true
}

\title{Chủ đề 8: Animation và cử chỉ nâng cao}
\subtitle{Lập trình Đa nền tảng -- Dart \& Flutter}
\author{Trương Nhật Hoàng}
\date{\today}

\begin{document}

% ======================================================
\begin{frame}
  \titlepage
\end{frame}

% ======================================================
\begin{frame}{1. Bài toán: thuộc tính UI thay đổi theo thời gian}
  Khi kích thước, màu, vị trí, trạng thái ẩn/hiện thay đổi, Flutter có 3 cách:
  \vspace{0.3cm}
  \begin{enumerate}
    \item \textbf{Implicit animation}: chỉ đổi giá trị, Flutter tự chạy chuyển động.
    \item \textbf{Explicit animation}: tự điều khiển bằng \texttt{AnimationController}
          (chạy, dừng, lặp, đảo chiều).
    \item \textbf{Gesture}: thao tác người dùng kích hoạt (1) hoặc (2).
  \end{enumerate}
  \vspace{0.3cm}
  Mức độ tự điều khiển tăng dần: $1 \rightarrow 2 \rightarrow 3$.
\end{frame}

% ======================================================
\begin{frame}[fragile]{2. Implicit: \texttt{AnimatedContainer}}
  \begin{itemize}
    \item Viết như \texttt{Container}, thêm \texttt{duration} và \texttt{curve}.
    \item \texttt{setState} đổi giá trị $\rightarrow$ tự chạy từ giá trị cũ sang mới.
  \end{itemize}
\begin{lstlisting}
GestureDetector(
  onTap: () => setState(() => _expanded = !_expanded),
  child: AnimatedContainer(
    duration: Duration(milliseconds: 300),
    curve: Curves.easeOut,
    height: _expanded ? 160 : 80,
    decoration: BoxDecoration(
      color: _expanded ? Colors.blue.shade100 : Colors.grey.shade200,
      borderRadius: BorderRadius.circular(_expanded ? 24 : 8),
    ),
  ),
)
\end{lstlisting}
  \begin{itemize}
    \item Đổi giá trị giữa chừng $\rightarrow$ chạy tiếp từ vị trí hiện tại.
  \end{itemize}
\end{frame}

% ======================================================
\begin{frame}[fragile]{3. Họ widget \texttt{Animated...} có sẵn}
  \begin{center}
  \begin{tabular}{ll}
    \toprule
    \textbf{Muốn đổi} & \textbf{Dùng} \\
    \midrule
    Kích thước, màu, bo góc, padding & \texttt{AnimatedContainer} \\
    Ẩn/hiện mờ dần & \texttt{AnimatedOpacity} \\
    Vị trí căn & \texttt{AnimatedAlign}, \texttt{AnimatedPositioned} \\
    Thay widget này bằng widget khác & \texttt{AnimatedSwitcher} \\
    Giá trị số tùy ý & \texttt{TweenAnimationBuilder} \\
    \bottomrule
  \end{tabular}
  \end{center}
\begin{lstlisting}
AnimatedSwitcher(
  duration: Duration(milliseconds: 200),
  child: Text(_expanded ? 'Chi tiet...' : 'Bam de mo rong',
              key: ValueKey(_expanded)),   // key khac -> widget moi
)
\end{lstlisting}
\end{frame}

% ======================================================
\begin{frame}{4. Khi implicit không đủ}
  Implicit chỉ chạy \textbf{một lần từ A sang B}. Không làm được:
  \vspace{0.2cm}
  \begin{itemize}
    \item Lặp vô hạn (icon loading xoay).
    \item Dừng giữa chừng, đảo chiều theo lệnh.
    \item Chạy theo vị trí tay kéo.
  \end{itemize}
  \vspace{0.3cm}
  $\Rightarrow$ Dùng \textbf{\texttt{AnimationController}}: một bộ đếm giá trị
  \texttt{value} từ 0.0 đến 1.0, cập nhật mỗi khung hình.
\end{frame}

% ======================================================
\begin{frame}[fragile]{5. \texttt{AnimationController} và quản lý tài nguyên}
  \begin{columns}[T]
  \begin{column}{0.55\textwidth}
\begin{lstlisting}
class _S extends State<DemoPage>
    with TickerProviderStateMixin {
  late final AnimationController _spin =
    AnimationController(
      vsync: this,
      duration: Duration(seconds: 1));

  @override
  void dispose() {
    _spin.dispose();   // bat buoc
    super.dispose();
  }
}
\end{lstlisting}
  \end{column}
  \begin{column}{0.43\textwidth}\small
    \begin{itemize}
      \item \texttt{vsync}: nhịp màn hình gọi cập nhật mỗi khung hình
            (tương tự ngắt timer).
      \item Tạo khi \texttt{State} khởi tạo.
      \item \texttt{dispose()} giải phóng (tương tự destructor C++).
      \item Quên \texttt{dispose()} $\rightarrow$ ticker vẫn chạy, rò rỉ tài nguyên.
      \item Mixin cung cấp \texttt{vsync}: \\
            {\scriptsize 1 controller: \texttt{SingleTicker\-ProviderStateMixin} \\
            nhiều controller: \texttt{TickerProviderStateMixin}}
    \end{itemize}
  \end{column}
  \end{columns}
\end{frame}

% ======================================================
\begin{frame}[fragile]{6. Gắn controller vào widget: họ \texttt{...Transition}}
  \begin{itemize}
    \item \texttt{ScaleTransition}, \texttt{RotationTransition},
          \texttt{FadeTransition}, \texttt{SlideTransition}.
    \item Tự vẽ: \texttt{AnimatedBuilder}.
    \item Lệnh: \texttt{forward()}, \texttt{reverse()}, \texttt{repeat()},
          \texttt{stop()}, \texttt{animateTo(x)}.
  \end{itemize}
\begin{lstlisting}
// Nut tim nay: 1.0 -> 1.4 -> 1.0
late final _heart = AnimationController(
    vsync: this, lowerBound: 1.0, upperBound: 1.4,
    duration: Duration(milliseconds: 150));

void _like() => _heart.forward(from: 1.0).then((_) => _heart.reverse());

IconButton(
  onPressed: _like,
  icon: ScaleTransition(scale: _heart,
      child: Icon(Icons.favorite, color: Colors.red)),
)

// Icon xoay lien tuc
RotationTransition(turns: _spin, child: Icon(Icons.sync))
onPressed: () => _spin.isAnimating ? _spin.stop() : _spin.repeat()
\end{lstlisting}
\end{frame}

% ======================================================
\begin{frame}[fragile]{7. Gesture: kích hoạt và điều khiển chuyển động}
  \begin{itemize}
    \item \texttt{GestureDetector}, \texttt{InkWell}: \texttt{onTap}, \texttt{onDoubleTap},
          \texttt{onLongPress}, \texttt{onHorizontalDragUpdate}, \dots
    \item Có sẵn hiệu ứng: \texttt{Dismissible} (vuốt xoá), \texttt{Draggable}, \texttt{Hero}.
  \end{itemize}
\begin{lstlisting}
GestureDetector(
  // Dang keo: tay dieu khien truc tiep value
  onHorizontalDragUpdate: (d) => _drag.value += d.primaryDelta! / 250,
  // Tha tay: controller tu chay ve mep gan nhat
  onHorizontalDragEnd: (_) => _drag.animateTo(_drag.value > 0.5 ? 1 : 0),
  child: AnimatedBuilder(
    animation: _drag,
    child: Container(width: 60, height: 60, color: Colors.teal),
    builder: (context, child) =>
        Align(alignment: Alignment(_drag.value * 2 - 1, 0), child: child),
  ),
)
\end{lstlisting}
\end{frame}

% ======================================================
\begin{frame}[fragile]{8. Bảo đảm chuyển động không cản trở thao tác}
  \begin{itemize}
    \item \textbf{Ngắn}: 150--300 ms cho phản hồi thao tác.
    \item \textbf{Không khoá input}: đang chạy vẫn bấm/kéo được;
          implicit tự chạy tiếp từ vị trí hiện tại.
    \item \textbf{Build ít}: phần tĩnh truyền qua tham số \texttt{child}
          của \texttt{AnimatedBuilder}.
    \item \textbf{Giảm chuyển động}: theo cài đặt hệ thống hoặc công tắc trong app.
  \end{itemize}
\begin{lstlisting}
bool get _reduce =>
    _reduceMotion || MediaQuery.of(context).disableAnimations;

Duration _ms(int ms) =>
    _reduce ? Duration.zero : Duration(milliseconds: ms);

AnimatedContainer(duration: _ms(300), ...)   // chuc nang giu nguyen,
                                              // chi bo chuyen dong
\end{lstlisting}
\end{frame}

% ======================================================
\begin{frame}{9. Demo: một màn hình, 4 nhóm hiệu ứng}
  \begin{center}\footnotesize
  \begin{tabular}{lll}
    \toprule
    \textbf{Thao tác} & \textbf{Hiệu ứng} & \textbf{Dùng} \\
    \midrule
    Bấm card & Mở rộng, đổi màu & \texttt{AnimatedContainer}, \texttt{AnimatedSwitcher} \\
    Bấm nút tim & Nảy lên rồi về & Controller + \texttt{ScaleTransition} \\
    Bấm nút xoay & Xoay liên tục / dừng & \texttt{repeat()}, \texttt{stop()} \\
    Kéo khối vuông & Theo tay, thả tự về mép & Gesture + \texttt{animateTo()} \\
    Vuốt task & Trượt ra và xoá & \texttt{Dismissible} \\
    Bật ``Giảm chuyển động'' & Mọi duration = 0 & \texttt{\_ms()} \\
    \bottomrule
  \end{tabular}
  \end{center}
  \vspace{0.2cm}
  3 controller đều được \texttt{dispose()} trong \texttt{dispose()} của \texttt{State}.
\end{frame}

% ======================================================
\begin{frame}{10. Tổng kết}
  \begin{itemize}
    \item Đổi giá trị A $\rightarrow$ B một lần: \textbf{implicit} (\texttt{Animated...}).
    \item Cần lặp, dừng, đảo chiều, theo tay: \textbf{\texttt{AnimationController}}
          + \texttt{...Transition}.
    \item Controller là tài nguyên: tạo trong \texttt{State}, \texttt{dispose()} khi huỷ.
    \item Gesture kích hoạt hoặc điều khiển trực tiếp \texttt{value}.
    \item Chuyển động ngắn, không khoá thao tác, tắt được khi cần.
  \end{itemize}
\end{frame}

\end{document}